\documentclass{oxmathproblems}

\printanswers

\course{Understanding Analysis, 2nd Edition}
\sheettitle{Chapter 6 - Section 6.2 - Uniform Convergence of a Sequence of Functions} 

\begin{document}
\begin{questions}

%------------------------- Problem 1 -------------------------
\miquestion Let
\[f_{n}(x)=\frac{nx}{1+nx^2}\]
\begin{enumerate}[label=(\alph*)]
  \item Find the pointwise limit of $(f_{n})$ for all $x \in (0, \infty)$.
  \item Is the convergence uniform on $(0, \infty)$?
  \item Is the convergence uniform on $(0, 1)$?
  \item Is the convergence uniform on $(1, \infty)$?
\end{enumerate}

\begin{solution}
  \begin{enumerate}[label=(\alph*)]
    \item Intuitively, as $n$ grows the $+ 1$ term starts to not matter, so it should converge to $\frac{nx}{nx^2}=\frac{1}{x}$. Formally
    \[\lim_{n \to \infty}\frac{nx}{1+nx^2}=\lim_{n \to \infty}\frac{x}{1/n + x^2}=1/x\]
    \item No. Given $\epsilon > 0$, let $N$ such that for all $n \geq N$, $|f_{n}(x)-\frac{1}{x}| < \epsilon$. Or in other words
    \[\left|\frac{nx}{1+nx^2}-\frac{1}{x}\right| = \frac{1}{x(1+nx^2)} < \epsilon\]
    For this to be uniformly convergent, the $N$ needs to work for all $x$. But the lhs is unbounded as $x \to 0$; there's always a small enough $x$
    such that the lhs is eventually greater than $\epsilon$.
    \item No, for the same reason as part (b). We can always choose small enough $x$ to break the uniform convergence.
    \item Now the answer is yes. Reusing the result from part (b), for $x > 1$, the value is upper-bounded by $x=1$
    \[\frac{1}{x(1+nx^2)} < \frac{1}{1+n}\]
    This doesn't depend on $x$, so we can simply choose $N \geq \frac{1}{\epsilon}-1$.
  \end{enumerate}
\end{solution}

%------------------------- Problem 2 -------------------------
\miquestion
\begin{enumerate}[label=(\alph*)]
  \item Define a sequence of functions on \textbf{R} by
  \[
    f_{n}(x) = \begin{cases}
      1 & \text{if } x = 1, \frac{1}{2}, \frac{1}{3}, ..., \frac{1}{n}\\
      0 & \text{otherwise}
    \end{cases}
  \]
  and let $f$ be the pointwise limit of $f_{n}$.

  Is each $f_{n}$ continuous at zero? Does $f_{n} \to f$ uniformly on \textbf{R}? Is $f$ continuous at zero?
  \item Repeat this exercise using the sequence of functions
    \[
      g_{n}(x) = \begin{cases}
        x & \text{if } x = 1, \frac{1}{2}, \frac{1}{3}, ..., \frac{1}{n}\\
        0 & \text{otherwise}
      \end{cases}
  \]
  \item Repeat once more with
    \[
      h_{n}(x) = \begin{cases}
        1 & \text{if } x = \frac{1}{n}\\
        x & \text{if } x = 1, \frac{1}{2}, \frac{1}{3}, ..., \frac{1}{n-1}\\
        0 & \text{otherwise}
      \end{cases}
  \]
\end{enumerate}
 In each case, explain how the results are consistent with the Continuous Limit Theorem (Theorem 6.2.6).

\begin{solution}
  \begin{enumerate}[label=(\alph*)]
    \item Yes, each $f_{n}$ is continuous at zero. Simply choose $\delta < \frac{1}{n}$, then $f(x)=0$ for $x < \delta$.
    
    To answer uniform convergence, let's first define $f$. For $x$ not of the form $\frac{1}{n}$, $f(x)$ is clearly 0. For $x$ of the form $\frac{1}{n}$, then
    for all $m \geq n$, $f_{m}(x)=1$ i.e. eventually all $f_{n}(x)$ becomes 1. So $f$ is
    \[
      f(x) = \begin{cases}
        1 & \text{if } x = \frac{1}{n} \text{ for } n \in \textbf{N}\\
        0 & \text{otherwise}
      \end{cases}
   \]

   It's not uniformly convergent. Let $\epsilon < 1$. Now say we choose $N=5$, then for $x = \frac{1}{6}$,
   we have $f_{5}(\frac{1}{6})=0$ but $f(\frac{1}{6})=1$.
   In general, no single $N$ can work for all $x$ because we can choose $x = \frac{1}{N+1}$, then $|f_{N}(x)-f(x)|=1 > \epsilon$.

   Finally, $f$ is not continuous at zero. We have $f(0)=0$. But since $\frac{1}{n} \to 0$, there's always $x$ of the form $\frac{1}{n}$ close enough to zero
   such that $f(x)=1$.

   \item Each $g_{n}$ is continuous at zero. Same argument as part (a), choose $\delta < \frac{1}{n}$ then $|x| < \delta \implies g_{n}(x)=0$.
   
   Using a very similar argument as part (a), the limit function $g$ is
    \[
      g(x) = \begin{cases}
        x & \text{if } x = \frac{1}{n} \text{ for } n \in \textbf{N}\\
        0 & \text{otherwise}
      \end{cases}
   \]
   
   But unlike part(a), this is uniformly convergent. The idea is to choose $N$ big enough such that even if $g_{N}(x)=0$, then $x$ needs to be really small, hence
   $|g_{N}(x)-g(x)|=|0-x|=|x|$ is really small as well.

   Formally, choose $N > \frac{1}{\epsilon}$. Then for $n \geq N$, either $g_{n}(x)=g(x)$ or if $g_{n}(x)=0$, then $x=\frac{1}{m}$ with $m > n$.
   But since $n > \frac{1}{\epsilon}$, if $m > n$ then $\frac{1}{m} < \epsilon$, hence $|g_{n}(\frac{1}{m})-g(\frac{1}{m})|=\frac{1}{m} < \epsilon$.

   $g$ is continuous at zero, as $\frac{1}{n}$ converges to zero.

   \item Each $h_{n}$ is continuous at zero. Choose $\delta < \frac{1}{n}$, then $|x| < \delta \implies h_{n}(x) = 0$.
   
   Using a similar argument as part (a), for each $x$ of the form $\frac{1}{n}$, we have $h_{m}(x)=x=\frac{1}{n}$ for all $m > n$. So $h$ is
    \[
      h(x) = \begin{cases}
        x & \text{if } x = \frac{1}{n} \text{ for } n \in \textbf{N}\\
        0 & \text{otherwise}
      \end{cases}
   \]
   
   $h$ is the same as $g$ from part(b), but in this case it's not uniformly convergent. The issue here is that
   \[h_{N}(\frac{1}{N})-h(\frac{1}{N})=1-\frac{1}{N}\]
   So as $x = \frac{1}{n}$ gets smaller, we need to choose increasingly larger $N$.

   More formally, let $\epsilon < \frac{1}{2}$ and $x = \frac{1}{2}$. Choose $N \geq 3$. Then for $n \geq N$, we have
   \[h_{n}(\frac{1}{2})-h(\frac{1}{2})=\frac{1}{2}-\frac{1}{2}=0\]
   But this same $N$ cannot for $x= \frac{1}{N+1}$ because for $n = N+1$ we have
   \[h_{N+1}(\frac{1}{N+1})-h(\frac{1}{N+1})=1-\frac{1}{N+1}\]
   And since $N \geq 3$, then $\frac{1}{N+1} \leq \frac{1}{4}$, hence
   \[1-\frac{1}{N+1} \geq 1-\frac{1}{4} > \epsilon\]

   Finally, $h$ is continuous at zero as it's the same as $g$ in part(b).
  \end{enumerate}
\end{solution}

%------------------------- Problem 3 -------------------------
\miquestion For each $n \in \textbf{N}$ and $x \in [0, \infty)$, let
\[g_{n}(x)=\frac{x}{1+x^{n}}\]
and 
  \[
    h_{n}(x) = \begin{cases}
      1 & \text{if } x \geq 1/n\\
      nx & \text{if } 0 \leq x < 1/n
    \end{cases}
  \]
\begin{enumerate}[label=(\alph*)]
  \item Find the pointwise limit on $[0, \infty)$.
  \item Explain how we know that the convergence \textit{cannot} be uniform on $[0, \infty)$.
  \item Choose a smaller set over which the convergence is uniform and supply a proof.
\end{enumerate}

\begin{solution}
  \begin{enumerate}[label=(\alph*)]
    \item
    \[
        g(x) = \begin{cases}
          x & \text{if } x < 1\\
          1/2 & \text{if } x = 1\\
          0 & \text{if } x > 1
        \end{cases}
    \]

    For $h$, note that for $x > 0$ and large enough $n$ then $x \geq 1/n$ eventually and we always enter the first case.\\
    For $x=0$, we always enter the second case. So $h$ is
    \[
        h(x) = \begin{cases}
          1 & \text{if } x > 0\\
          0 & \text{if } x = 0
        \end{cases}
    \]

    \item For $h_{n}$, let's try $\epsilon=1/4$ and $x=1/2$ and see what we get.
    
    $h_{1}(1/2)=\frac{1}{2}$ and $|h_{1}(1/2)-h(1/2)|=|1/2 - 1|=1/2 > \epsilon$, so $N=1$ doesn't work.

    $h_{2}(1/2)=1$. So $N \geq 2$ works because $h_{N}(1/2)=1$.

    But this $N$ doesn't work for $x = \frac{1}{N^{2}}$, because $h_{N}(\frac{1}{N^2})=\frac{1}{N}$ and even if we choose the minimum $N=2$, we
    get $|\frac{1}{2}-1|=\frac{1}{2} > \epsilon$. Choosing larger $N$ will only increase the distance from 1 even more.

    \textbf{-- Better solution by chatgpt --}\\
    Use the Continuous Limit Theorem. $g_{n}$ and $h_{n}$ is continuous for all $n$, so if they're uniformly convergent then $g$ and $h$ must be continuous.
    But $g$ is discontinuous at $x=1$ and $h$ at $x=0$.

    \item For $g_{n}$, a simple option is to restrict to $x=1$. Then $g_{n}(1)=g(1)=1/2$ for all $n$ which is certainly uniformly convergent. Likewise for $h_{n}$,
    restrict to $x=0$. Then $h_{n}(0)=h(0)=0$ for all $n$. But such trivial sets are probably not what the author intended. Let's choose a bigger set.

    \textbf{-- Incorrect -- }\\
    For $g_{n}$, we can stay away from the discontinuity at $1$. Say, $(1, \infty)$. Then $g(x)=0$ and
    \[|g_{n}(x)-g(x)|=\frac{x}{1+x^{n}} \leq \frac{x}{x^n}=\frac{1}{x^{n-1}}\]

    \textbf{-- Corrected --}\\
    Oops, $(1, \infty)$ doesn't work because $x$ can still get arbitrarily close to 1. For any $N$, we can choose $x$ close enough to 1
    such that $\frac{1}{x^{n-1}}$ gets close to 1 hence far from $g(x)=0$.

    Instead, we need an interval bounded away from 1, say $[2, \infty)$. then
    \[|g_{n}(x)-g(x)| \leq \frac{1}{x^{n-1}} \leq \frac{1}{2^{n-1}} \to 0\]
    So for any $\epsilon$, choose $N$ big enough so that $\frac{1}{2^{N-1}} < \epsilon$.

    $h_{n}$ is uniformly convergent on $[1, \infty)$ because for $x \geq 1$ we'll always enter the first case, hence $h_{n}(x)=1$ for all $n$ and $x$.
  \end{enumerate}
\end{solution}

%------------------------- Problem 4 -------------------------
\miquestion TODO. It depends on Exercise 5.2.8 which I haven't done yet.

%------------------------- Problem 5 -------------------------
\miquestion Using the Cauchy Criterion for convergent sequences of real numbers (Theorem 2.6.4), supply a proof for Theorem 6.2.5. (First, define a candidate for $f(x)$,
and then argue that $f_{n} \to f$ uniformly).
\begin{solution}
  The proof strategy is very similar to Theorem 2.6.4.

  \textbf{Necessity: uniform convergence $\implies$ Cauchy}\\
  There exists $N$ such that for all $n \geq N$, $|f_{n}(x)-f(x)| < \epsilon/2$ for all $x$. Then for any $m, n \geq N$ and any $x$ we have
  \[|f_{n}(x)-f_{m}(x)| \leq |f_{n}(x) - f(x)| + |f(x)-f_{m}(x)| < \epsilon/2 + \epsilon/2 = \epsilon\]
  Satisfying Cauchy criterion.

  \textbf{Sufficiency: Cauchy $\implies$ uniform convergence}\\
  The Cauchy criterion doesn't say what the limit function is, so we need to define it first. For a \textbf{particular} $x$, the sequence of real numbers
  $f_{1}(x), f_{2}(x), ...$ is a Cauchy sequence, so it converges to a limit $L_{x}$. Define $f$ to be
  \[f(x)=L_{x}\]

  Now we prove uniform convergence. By the Cauchy hypothesis (for sequence of functions), there exists $N$ such that for all $n, m \geq N$, we have
  \[|f_{n}(x)-f_{m}(x)| < \epsilon/2\]
  for any $x$.

  Also, by the fact that $f_{1}(x), f_{2}(x), ... \to L_{x}$, there exists a $k \geq N$ such that $|f_{k}(x)-L_{x}| < \epsilon/2$.
  The value of $k$ depends on $x$, but that's okay because for any $x$, $k \geq N$ and $N$ does not depend on $x$.

  Combining these two inequalities, we satisfy the condition for uniform convergence. For any $n \geq N$, we have
  \[|f_{n}(x)-L_{x}| \leq |f_{n}(x)-f_{k}(x)| + |f_{k}(x)-L_{x}| < \epsilon/2 + \epsilon/2 = \epsilon\]
  for any $x$.
\end{solution}

%------------------------- Problem 6 -------------------------
\miquestion TODO

%------------------------- Problem 7 -------------------------
\miquestion Let $f$ be uniformly continuous on all of \textbf{R}, and define a sequence of functions by $f_{n}(x)=f(x+\frac{1}{n})$. Show that $f_{n} \to f$ uniformly.
Give an example to show that this proposition fails if $f$ is only assumed to be continuous and not uniformly continuous on \textbf{R}.
\begin{solution}
  We need to show there exists $N$ such that for all $n \geq N$ and for all $x$
  \[\left|f(x+\frac{1}{n})-f(x)\right| < \epsilon\]

  By uniform continuity, there exists $\delta > 0$ such that for all $x$, if $x-\delta < c < x+\delta$ then $|f(c)-f(x)| < \epsilon$.
  The idea is we just have to pick $N$ (for the uniform convergence) large enough such that $\frac{1}{N} < \delta$.

  Since $(\frac{1}{n}) \to 0$, choose $N$ large enough such that $\frac{1}{N} < \delta$. Then for all $n \geq N$
  \[x-\delta < x+\frac{1}{n} < x + \delta \implies \left|f(x+\frac{1}{n})-f(x)\right| < \epsilon\]

  For the counterexample, take $f(x)=x^2$. This is continuous but not uniformly continuous on \textbf{R}. Then
  \[\left|f(x+\frac{1}{n})-f(x)\right| = \left|\frac{1}{n^2} + \frac{2x}{n}\right|\]

  Suppose there's an $N$ that works for a particular $x_{1}$ and $\epsilon$. There exists another $x_{2}$ large enough such that $\frac{2x}{n}$ is eventually
  larger than $\epsilon$, so the convergence is not uniform. More formally, let $\epsilon = 1$. If we choose $x=N$ then
  \[\left|\frac{1}{n^2} + \frac{2x}{n}\right| = 2 + \frac{1}{n^2} > 1\]

  As to how I knew that $x^2$ would work as a counterexample, I don't. I just chose the simplest non uniformly continuous function I can think of and it luckily worked.
  Chatgpt told me the intuition behind this is that "$x^2$ changes faster if x is large, so shifting by $1/n$ may be tiny, but at large $x$ the vertical change can still be large".
\end{solution}

%------------------------- Problem 8 -------------------------
\miquestion Let $(g_{n})$ be a sequence of continuous functions that converges uniformly to $g$ on a compact set $K$. If $g(x) \neq 0$ on $K$, show $(1/g_{n})$ converges
uniformly on $K$ to $1/g$.

\begin{solution}
  I missed the fact that we can use compactness to bound $g(K)$ and only realized it with chatgpt's assistance.

  Given $\epsilon > 0$, we need to show that there exists $N$ such that for all $n \geq N$ and for all $x \in K$
  \[\left|\frac{1}{g_{n}(x)}-\frac{1}{g(x)}\right| = \frac{|g_{n}(x)-g(x)|}{|g_{n}(x)||g(x)|} < \epsilon\]

  The difficulty lies in bounding the denominator. We use the fact that $g$ is continuous on a compact set. By Theorem 4.4.1, $g(K)$ is compact, so it's bounded.
  And since $0 \notin g(K)$, then there exists a positive lower bound $m$ such that $|g(x)| \geq m$, then
  \[\frac{|g_{n}(x)-g(x)|}{|g_{n}(x)||g(x)|} \leq \frac{|g_{n}(x)-g(x)|}{|g_{n}(x)|m}\]
  
  Since $(g_{n})$ converges uniformly to $g$, we can choose $N_{1}$ such that $g_{n}(x)$ is close to $g(x)$ or equivalently, $|g_{n}(x)|$ close to $|g(x)|$. Formally,
  this follows from the fact that
  \[||g_{n}(x)|-|g(x)|| \leq |g_{n}(x)-g(x)|\]

   Choose $N_{1}$ such that $|g_{n}(x)| \geq \frac{m}{2}$, resulting in
  \[\frac{|g_{n}(x)-g(x)|}{|g_{n}(x)|m} \leq \frac{2}{m^2}|g_{n}(x)-g(x)|\]

  Finally, choose $N_{2}$ such that $|g_{n}(x)-g(x)| < \frac{\epsilon m^{2}}{2}$ and choose $N=\max(N_{1}, N_{2})$
  \[\frac{2}{m^2}|g_{n}(x)-g(x)| < \frac{2}{m^2} \frac{\epsilon m^{2}}{2} = \epsilon\]

\end{solution}

\end{questions}
\end{document}
